\documentclass[11pt,a4paper]{article}
\usepackage[utf8]{inputenc}
\usepackage[T1]{fontenc}
\usepackage[spanish,es-minimal,es-noshorthands]{babel}
\usepackage{lmodern}
\usepackage[sfdefault,lining]{FiraSans}
\usepackage[a4paper,top=2.3cm,bottom=2.2cm,left=1.8cm,right=1.8cm,headheight=1.3cm,headsep=0.4cm,footskip=1cm]{geometry}
\usepackage{amsmath,amssymb}
\allowdisplaybreaks
\usepackage{icomma}
\usepackage{graphicx}
\usepackage[table]{xcolor}
\usepackage{tikz}
\usepackage{fontawesome5}
\usepackage{ccicons}
\usepackage{fancyhdr}
\usepackage{lastpage}
\usepackage{tasks}
\usepackage{needspace}
\usepackage[most]{tcolorbox}
\usepackage{array}
\usepackage[hidelinks]{hyperref}
\AtBeginDocument{\hypersetup{pdftitle={Ficha de ejercicios de la unidad 6: Cálculo de derivadas. Derivabilidad de una función. 2º Bach.},pdfsubject={Matemáticas II 2º Bach.},pdfkeywords={derivada, recta tangente, recta normal, regla de la cadena, derivación logarítmica, derivación implícita, derivabilidad, funciones a trozos, parámetros, PAU}}}

\definecolor{navy}{HTML}{1B2748}
\definecolor{accent}{HTML}{E8603A}
\definecolor{teal}{HTML}{1F9E90}
\definecolor{soft}{HTML}{EEF1F8}
\definecolor{muted}{HTML}{5D6685}

\setlength{\parindent}{0pt}
\setlength{\parskip}{3pt}
\settasks{label=\alph*), label-format=\bfseries\color{navy}, label-width=1.3em, item-indent=1.8em, column-sep=0.6em, after-item-skip=5pt, before-skip=3pt, after-skip=4pt}

% ---------- capçalera i peu ----------
\newcommand{\llicencia}{\href{https://creativecommons.org/licenses/by-nc-sa/4.0/deed.es}{\ccbyncsa}\ \ Licencia Creative Commons Atribución-NoComercial-CompartirIgual 4.0 Internacional (CC BY-NC-SA 4.0)}
\pagestyle{fancy}
\fancyhf{}
\fancyhead[L]{\raisebox{-0.35\height}{\includegraphics[height=0.95cm]{logo-gva.png}}}
\fancyhead[C]{\small\color{muted}Ficha de ejercicios · Unidad 6: Cálculo de derivadas · 2\textsuperscript{o} Bach.}
\fancyhead[R]{\raisebox{-0.35\height}{\includegraphics[height=0.95cm]{logo-vermellar.png}}}
\fancyfoot[L]{\scriptsize\color{muted}\llicencia}
\fancyfoot[R]{\small\color{muted}\thepage}
\renewcommand{\headrulewidth}{0.4pt}
\renewcommand{\headrule}{\color{muted!40}\hrule height \headrulewidth}
\fancypagestyle{primera}{\fancyhf{}\renewcommand{\headrulewidth}{0pt}
  \fancyfoot[L]{\scriptsize\color{muted}\llicencia}
  \fancyfoot[R]{\small\color{muted}\thepage}}

% ---------- nivells i exercicis ----------
\newcommand{\estrelles}[1]{{\color{accent}\ifcase#1\or\faStar\faStar[regular]\faStar[regular]\or\faStar\faStar\faStar[regular]\or\faStar\faStar\faStar\fi}}
\newcounter{exer}
\newcommand{\ex}[1]{\par\needspace{4\baselineskip}\medskip\refstepcounter{exer}\noindent{\color{navy}\bfseries\theexer.}\ \ }

\newcommand{\seccio}[2]{\par\needspace{10\baselineskip}\bigskip
  \begin{tcolorbox}[colback=navy,colframe=navy,coltext=white,arc=2mm,boxrule=0pt,left=3mm,right=3mm,top=1.2mm,bottom=1.2mm]
  {\large\bfseries #1\quad #2}
  \end{tcolorbox}\vspace{-2pt}}


\setlength{\emergencystretch}{3em}

% quadre de recordatori (només a la fitxa d'exercicis)
\newcommand{\recorda}[1]{\par\smallskip\begin{tcolorbox}[colback=accent!7,colframe=accent!55,arc=2mm,boxrule=0.5pt,left=2.5mm,right=2.5mm,top=1mm,bottom=1mm]\small{\color{accent}\faLightbulb[regular]}\ \textbf{Recuerda:}\ #1\end{tcolorbox}}
% etiqueta dels exercicis de la PAU d'altres comunitats
\newcommand{\pau}[1]{{\color{teal}\small\bfseries\faUniversity\ PAU · #1.}\ }
\newcommand{\res}[1]{#1}\newcommand{\rest}[1]{#1}

% =====================================================================
% Capçalera FentMates de les fitxes: un color per a cada curs.
% Els colors són els de la pàgina inicial (LEVELS a assets/js/materials.js);
% el logo ve de _eines/logo/pdf/ (python3 _eines/logo/genera.py).
%
% Ús, al final del preàmbul de la unitat:
%   % =====================================================================
% Capçalera FentMates de les fitxes: un color per a cada curs.
% Els colors són els de la pàgina inicial (LEVELS a assets/js/materials.js);
% el logo ve de _eines/logo/pdf/ (python3 _eines/logo/genera.py).
%
% Ús, al final del preàmbul de la unitat:
%   % =====================================================================
% Capçalera FentMates de les fitxes: un color per a cada curs.
% Els colors són els de la pàgina inicial (LEVELS a assets/js/materials.js);
% el logo ve de _eines/logo/pdf/ (python3 _eines/logo/genera.py).
%
% Ús, al final del preàmbul de la unitat:
%   \input{../comu/fentmates.tex}
%   \fmcurs{1eso}
% i, al començament del document:
%   \capcalera{Matemàtiques}{1r ESO}{Fitxa d'exercicis · Unitat 2}{Divisibilitat}
% Opcional abans de \input: \def\fmneed{14} (espai mínim abans de cada secció).
% =====================================================================
\makeatletter
\def\fm@logodir{../../logo/pdf}
\providecommand{\fmneed}{10}
\def\fm@cursnom{}

% clau del curs, color i color del text damunt del color
\newcommand{\fm@curs}[3]{\@namedef{fm@color@#1}{#2}\@namedef{fm@text@#1}{#3}}
\fm@curs{1eso}{E6B84A}{navy}
\fm@curs{2eso}{3B82F6}{white}
\fm@curs{3eso}{FF7A50}{white}
\fm@curs{4eso-a}{A78BFA}{white}
\fm@curs{4eso-b}{6D28D9}{white}
\fm@curs{1batx-mi}{9B1C1C}{white}
\fm@curs{1batx-ccss}{EF6B6B}{white}
\fm@curs{2batx-mii}{1E6B47}{white}
\fm@curs{2batx-ccss}{2BB3A3}{white}

\newcommand{\fmcurs}[1]{%
  \@ifundefined{fm@color@#1}{\PackageError{fentmates}{Curs desconegut: #1}{}}{}%
  \edef\fm@tmp{\noexpand\definecolor{curs}{HTML}{\@nameuse{fm@color@#1}}%
               \noexpand\colorlet{curstext}{\@nameuse{fm@text@#1}}}\fm@tmp
  \colorlet{cursfosc}{curs!72!black}%
  \def\fm@logo{\fm@logodir/vertical-#1.pdf}%
  \def\fm@marca{\fm@logodir/marca-#1.pdf}}

% ---------- capçalera de la primera pàgina ----------
% #1 matèria (petit), #2 curs, #3 tipus de document i unitat, #4 títol
\newcommand{\capcalera}[4]{%
  \gdef\fm@cursnom{#2}%
  \par\noindent
  \begin{tcolorbox}[colback=curs,colframe=curs,boxrule=0pt,arc=2.2mm,width=3.5cm,height=2.15cm,
      valign=center,left=3mm,right=2mm,top=1mm,bottom=1mm,nobeforeafter,box align=center]
    \color{curstext}{\footnotesize #1}\par\vspace{1pt}{\fontsize{19}{21}\selectfont\bfseries #2\par}%
  \end{tcolorbox}\hspace{5mm}%
  \begin{minipage}[c]{\dimexpr\linewidth-3.5cm-5mm-2.9cm\relax}
    {\color{cursfosc}\small\bfseries\MakeUppercase{#3}}\par\vspace{2pt}
    {\color{navy}\fontsize{20}{23}\selectfont\bfseries #4\par}
  \end{minipage}\hfill
  \begin{minipage}[c]{2.6cm}\raggedleft\includegraphics[width=2.6cm]{\fm@logo}\end{minipage}\par
  \vspace{2pt}}

% ---------- barres de secció del color del curs ----------
\renewcommand{\seccio}[2]{\par\needspace{\fmneed\baselineskip}\bigskip
  \begin{tcolorbox}[colback=curs,colframe=curs,coltext=curstext,arc=2mm,boxrule=0pt,left=3mm,right=3mm,top=1.2mm,bottom=1.2mm]
  {\large\bfseries #1\quad #2}
  \end{tcolorbox}\vspace{-2pt}}

% ---------- estrelles de dificultat del color del curs (a Batxillerat no se'n posen) ----------
\@ifundefined{estrelles}{}{%
  \renewcommand{\estrelles}[1]{{\color{curs}\ifcase#1\or\faStar\faStar[regular]\faStar[regular]\or\faStar\faStar\faStar[regular]\or\faStar\faStar\faStar\fi}}}

% ---------- capçalera i peu de les pàgines ----------
% primera pàgina: logos institucionals al peu
\fancypagestyle{primera}{\fancyhf{}\renewcommand{\headrulewidth}{0pt}
  \fancyfoot[L]{\raisebox{-0.32\height}{\includegraphics[height=0.78cm]{logo-gva.png}}\hspace{5mm}%
    \raisebox{-0.32\height}{\includegraphics[height=0.78cm]{logo-vermellar.png}}}
  \fancyfoot[R]{\scriptsize\color{muted}\href{https://creativecommons.org/licenses/by-nc-sa/4.0/deed.ca}{\ccbyncsa}\ \ CC BY-NC-SA 4.0 · fentmates.org\qquad\small\thepage}}
% la resta: el símbol a l'esquerra i el curs a la dreta, amb el títol al centre com abans
\fancyhead[L]{\raisebox{-0.3\height}{\includegraphics[height=0.62cm]{\fm@marca}}}
\fancyhead[R]{\small\bfseries\color{cursfosc}\fm@cursnom}
\renewcommand{\headrule}{\color{curs!45}\hrule height \headrulewidth}
\makeatother

%   \fmcurs{1eso}
% i, al començament del document:
%   \capcalera{Matemàtiques}{1r ESO}{Fitxa d'exercicis · Unitat 2}{Divisibilitat}
% Opcional abans de \input: \def\fmneed{14} (espai mínim abans de cada secció).
% =====================================================================
\makeatletter
\def\fm@logodir{../../logo/pdf}
\providecommand{\fmneed}{10}
\def\fm@cursnom{}

% clau del curs, color i color del text damunt del color
\newcommand{\fm@curs}[3]{\@namedef{fm@color@#1}{#2}\@namedef{fm@text@#1}{#3}}
\fm@curs{1eso}{E6B84A}{navy}
\fm@curs{2eso}{3B82F6}{white}
\fm@curs{3eso}{FF7A50}{white}
\fm@curs{4eso-a}{A78BFA}{white}
\fm@curs{4eso-b}{6D28D9}{white}
\fm@curs{1batx-mi}{9B1C1C}{white}
\fm@curs{1batx-ccss}{EF6B6B}{white}
\fm@curs{2batx-mii}{1E6B47}{white}
\fm@curs{2batx-ccss}{2BB3A3}{white}

\newcommand{\fmcurs}[1]{%
  \@ifundefined{fm@color@#1}{\PackageError{fentmates}{Curs desconegut: #1}{}}{}%
  \edef\fm@tmp{\noexpand\definecolor{curs}{HTML}{\@nameuse{fm@color@#1}}%
               \noexpand\colorlet{curstext}{\@nameuse{fm@text@#1}}}\fm@tmp
  \colorlet{cursfosc}{curs!72!black}%
  \def\fm@logo{\fm@logodir/vertical-#1.pdf}%
  \def\fm@marca{\fm@logodir/marca-#1.pdf}}

% ---------- capçalera de la primera pàgina ----------
% #1 matèria (petit), #2 curs, #3 tipus de document i unitat, #4 títol
\newcommand{\capcalera}[4]{%
  \gdef\fm@cursnom{#2}%
  \par\noindent
  \begin{tcolorbox}[colback=curs,colframe=curs,boxrule=0pt,arc=2.2mm,width=3.5cm,height=2.15cm,
      valign=center,left=3mm,right=2mm,top=1mm,bottom=1mm,nobeforeafter,box align=center]
    \color{curstext}{\footnotesize #1}\par\vspace{1pt}{\fontsize{19}{21}\selectfont\bfseries #2\par}%
  \end{tcolorbox}\hspace{5mm}%
  \begin{minipage}[c]{\dimexpr\linewidth-3.5cm-5mm-2.9cm\relax}
    {\color{cursfosc}\small\bfseries\MakeUppercase{#3}}\par\vspace{2pt}
    {\color{navy}\fontsize{20}{23}\selectfont\bfseries #4\par}
  \end{minipage}\hfill
  \begin{minipage}[c]{2.6cm}\raggedleft\includegraphics[width=2.6cm]{\fm@logo}\end{minipage}\par
  \vspace{2pt}}

% ---------- barres de secció del color del curs ----------
\renewcommand{\seccio}[2]{\par\needspace{\fmneed\baselineskip}\bigskip
  \begin{tcolorbox}[colback=curs,colframe=curs,coltext=curstext,arc=2mm,boxrule=0pt,left=3mm,right=3mm,top=1.2mm,bottom=1.2mm]
  {\large\bfseries #1\quad #2}
  \end{tcolorbox}\vspace{-2pt}}

% ---------- estrelles de dificultat del color del curs (a Batxillerat no se'n posen) ----------
\@ifundefined{estrelles}{}{%
  \renewcommand{\estrelles}[1]{{\color{curs}\ifcase#1\or\faStar\faStar[regular]\faStar[regular]\or\faStar\faStar\faStar[regular]\or\faStar\faStar\faStar\fi}}}

% ---------- capçalera i peu de les pàgines ----------
% primera pàgina: logos institucionals al peu
\fancypagestyle{primera}{\fancyhf{}\renewcommand{\headrulewidth}{0pt}
  \fancyfoot[L]{\raisebox{-0.32\height}{\includegraphics[height=0.78cm]{logo-gva.png}}\hspace{5mm}%
    \raisebox{-0.32\height}{\includegraphics[height=0.78cm]{logo-vermellar.png}}}
  \fancyfoot[R]{\scriptsize\color{muted}\href{https://creativecommons.org/licenses/by-nc-sa/4.0/deed.ca}{\ccbyncsa}\ \ CC BY-NC-SA 4.0 · fentmates.org\qquad\small\thepage}}
% la resta: el símbol a l'esquerra i el curs a la dreta, amb el títol al centre com abans
\fancyhead[L]{\raisebox{-0.3\height}{\includegraphics[height=0.62cm]{\fm@marca}}}
\fancyhead[R]{\small\bfseries\color{cursfosc}\fm@cursnom}
\renewcommand{\headrule}{\color{curs!45}\hrule height \headrulewidth}
\makeatother

%   \fmcurs{1eso}
% i, al començament del document:
%   \capcalera{Matemàtiques}{1r ESO}{Fitxa d'exercicis · Unitat 2}{Divisibilitat}
% Opcional abans de \input: \def\fmneed{14} (espai mínim abans de cada secció).
% =====================================================================
\makeatletter
\def\fm@logodir{../../logo/pdf}
\providecommand{\fmneed}{10}
\def\fm@cursnom{}

% clau del curs, color i color del text damunt del color
\newcommand{\fm@curs}[3]{\@namedef{fm@color@#1}{#2}\@namedef{fm@text@#1}{#3}}
\fm@curs{1eso}{E6B84A}{navy}
\fm@curs{2eso}{3B82F6}{white}
\fm@curs{3eso}{FF7A50}{white}
\fm@curs{4eso-a}{A78BFA}{white}
\fm@curs{4eso-b}{6D28D9}{white}
\fm@curs{1batx-mi}{9B1C1C}{white}
\fm@curs{1batx-ccss}{EF6B6B}{white}
\fm@curs{2batx-mii}{1E6B47}{white}
\fm@curs{2batx-ccss}{2BB3A3}{white}

\newcommand{\fmcurs}[1]{%
  \@ifundefined{fm@color@#1}{\PackageError{fentmates}{Curs desconegut: #1}{}}{}%
  \edef\fm@tmp{\noexpand\definecolor{curs}{HTML}{\@nameuse{fm@color@#1}}%
               \noexpand\colorlet{curstext}{\@nameuse{fm@text@#1}}}\fm@tmp
  \colorlet{cursfosc}{curs!72!black}%
  \def\fm@logo{\fm@logodir/vertical-#1.pdf}%
  \def\fm@marca{\fm@logodir/marca-#1.pdf}}

% ---------- capçalera de la primera pàgina ----------
% #1 matèria (petit), #2 curs, #3 tipus de document i unitat, #4 títol
\newcommand{\capcalera}[4]{%
  \gdef\fm@cursnom{#2}%
  \par\noindent
  \begin{tcolorbox}[colback=curs,colframe=curs,boxrule=0pt,arc=2.2mm,width=3.5cm,height=2.15cm,
      valign=center,left=3mm,right=2mm,top=1mm,bottom=1mm,nobeforeafter,box align=center]
    \color{curstext}{\footnotesize #1}\par\vspace{1pt}{\fontsize{19}{21}\selectfont\bfseries #2\par}%
  \end{tcolorbox}\hspace{5mm}%
  \begin{minipage}[c]{\dimexpr\linewidth-3.5cm-5mm-2.9cm\relax}
    {\color{cursfosc}\small\bfseries\MakeUppercase{#3}}\par\vspace{2pt}
    {\color{navy}\fontsize{20}{23}\selectfont\bfseries #4\par}
  \end{minipage}\hfill
  \begin{minipage}[c]{2.6cm}\raggedleft\includegraphics[width=2.6cm]{\fm@logo}\end{minipage}\par
  \vspace{2pt}}

% ---------- barres de secció del color del curs ----------
\renewcommand{\seccio}[2]{\par\needspace{\fmneed\baselineskip}\bigskip
  \begin{tcolorbox}[colback=curs,colframe=curs,coltext=curstext,arc=2mm,boxrule=0pt,left=3mm,right=3mm,top=1.2mm,bottom=1.2mm]
  {\large\bfseries #1\quad #2}
  \end{tcolorbox}\vspace{-2pt}}

% ---------- estrelles de dificultat del color del curs (a Batxillerat no se'n posen) ----------
\@ifundefined{estrelles}{}{%
  \renewcommand{\estrelles}[1]{{\color{curs}\ifcase#1\or\faStar\faStar[regular]\faStar[regular]\or\faStar\faStar\faStar[regular]\or\faStar\faStar\faStar\fi}}}

% ---------- capçalera i peu de les pàgines ----------
% primera pàgina: logos institucionals al peu
\fancypagestyle{primera}{\fancyhf{}\renewcommand{\headrulewidth}{0pt}
  \fancyfoot[L]{\raisebox{-0.32\height}{\includegraphics[height=0.78cm]{logo-gva.png}}\hspace{5mm}%
    \raisebox{-0.32\height}{\includegraphics[height=0.78cm]{logo-vermellar.png}}}
  \fancyfoot[R]{\scriptsize\color{muted}\href{https://creativecommons.org/licenses/by-nc-sa/4.0/deed.ca}{\ccbyncsa}\ \ CC BY-NC-SA 4.0 · fentmates.org\qquad\small\thepage}}
% la resta: el símbol a l'esquerra i el curs a la dreta, amb el títol al centre com abans
\fancyhead[L]{\raisebox{-0.3\height}{\includegraphics[height=0.62cm]{\fm@marca}}}
\fancyhead[R]{\small\bfseries\color{cursfosc}\fm@cursnom}
\renewcommand{\headrule}{\color{curs!45}\hrule height \headrulewidth}
\makeatother

\fmcurs{2batx-mii}

\begin{document}
\thispagestyle{primera}

\capcalera{Matemáticas II}{2\textsuperscript{o} Bach.}{Ficha de ejercicios · Unidad 6}{\raggedright Cálculo de derivadas. Derivabilidad de una función}
\vspace{10pt}
Nombre y apellidos: \rule{7.2cm}{0.4pt}\hfill Grupo: \rule{1.6cm}{0.4pt}\hfill Fecha: \rule{2.2cm}{0.4pt}
\vspace{8pt}

\begin{tcolorbox}[colback=soft,colframe=soft,arc=2mm,left=3mm,right=3mm,top=1.5mm,bottom=1.5mm]
\small Resuelve los ejercicios en el cuaderno justificando cada paso. Los ejercicios marcados con {\color{teal}\bfseries\faUniversity\ PAU} son de las pruebas de acceso a la universidad de otras comunidades autónomas (de Matemáticas II y de Matemáticas Aplicadas a las Ciencias Sociales II), traducidos al castellano; a los marcados como «ampliado» se les ha añadido el estudio de la derivabilidad. Puedes comprobar los resultados con las herramientas de los apuntes interactivos.
\end{tcolorbox}

% =====================================================
\seccio{1}{Tasa de variación y derivada por definición}

\ex{1} Calcula la tasa de variación media de las funciones siguientes en el intervalo indicado:\begin{tasks}(3)
  \task $f(x)=x^3-x$ en $[0,2]$
  \task $g(x)=\ln x$ en $[1,e]$
  \task $h(x)=\sin x$ en $\left[0,\frac\pi2\right]$
\end{tasks}\recorda{$\text{TVM}[a,b]=\dfrac{f(b)-f(a)}{b-a}$: es la pendiente de la recta secante. La derivada es el límite de la TVM cuando el intervalo se hace tan pequeño como queramos.}

\ex{2} Lanzamos una pelota hacia arriba y su altura, en metros, a los $t$ segundos es $s(t)=-5t^2+20t$. Calcula la velocidad media en los dos primeros segundos y, aplicando la definición de derivada, la velocidad instantánea en $t=1$.

\ex{2} Calcula, aplicando la definición de derivada:\begin{tasks}(2)
  \task $f'(2)$ si $f(x)=\dfrac{1}{x^2}$
  \task $f'(1)$ si $f(x)=\dfrac{x}{x+1}$
  \task $f'(0)$ si $f(x)=e^{3x}$
  \task $f'\!\left(\frac\pi2\right)$ si $f(x)=\sin x$
\end{tasks}\recorda{$f'(a)=\displaystyle\lim_{h\to0}\frac{f(a+h)-f(a)}{h}$. Con exponenciales y funciones trigonométricas, usa los infinitésimos equivalentes: $e^{u}-1\sim u$, \ $1-\cos u\sim\frac{u^2}{2}$, \ $\sin u\sim u$.}

\ex{3} Calcula la función derivada aplicando la definición:\begin{tasks}(2)
  \task $f(x)=x^3$
  \task $f(x)=\dfrac{1}{\sqrt x}$
  \task $f(x)=\cos x$
  \task $f(x)=\ln(2x)$
\end{tasks}

\ex{2} \pau{Navarra, extraordinaria 2023} Considera la función $g(x)=-x^2+4x+3$. Calcula la derivada de $g(x)$ en el punto $x=1$ aplicando la definición de derivada.


% =====================================================
\seccio{2}{Rectas tangente y normal}

\ex{1} Escribe las ecuaciones de las rectas tangente y normal a la gráfica de cada función en el punto indicado:\begin{tasks}(2)
  \task $f(x)=x^3-3x+1$ en $x=2$
  \task $g(x)=\ln(x^2-3)$ en $x=2$
\end{tasks}\recorda{Tangente: $y-f(a)=f'(a)(x-a)$. \ Normal (perpendicular a la tangente): $y-f(a)=-\dfrac{1}{f'(a)}(x-a)$; si $f'(a)=0$, la normal es la recta vertical $x=a$.}

\ex{2} Halla la ecuación de la recta tangente a la gráfica de cada función en el punto indicado:\begin{tasks}(1)
  \task \pau{Murcia, julio 2020} $f(x)=\dfrac{1}{x+1}$ en $x=1$.
  \task \pau{Navarra, ordinaria 2023} $f(x)=\sqrt{5+x^2}$ en $x=-2$.
  \task \pau{Navarra, ordinaria 2022} $f(x)=\dfrac{x}{x^2-4}$ en $x=1$.
\end{tasks}

\ex{2} \pau{Islas Baleares, extraordinaria 2021} Considera la función $f(x)=\dfrac3x+8$.\begin{tasks}(1)
  \task Halla los puntos de la gráfica en los que la recta tangente es paralela a la recta $3x+4y+5=0$.
  \task Calcula las ecuaciones de las rectas tangentes en esos puntos.
\end{tasks}

\ex{2} \pau{Murcia, septiembre 2018} Dada la función $f(x)=5x^3e^{2x}+\dfrac{1}{x^2+1}$, calcula $f'(0)$ y halla la ecuación de la recta tangente a la gráfica en el punto $(0,1)$.

\ex{2} Dada $f(x)=x^3-6x^2+9x$, halla los puntos donde la tangente es horizontal y la ecuación de la recta tangente paralela a $y=-3x$.

\ex{3} Halla las rectas tangentes a la parábola $y=x^2$ que pasan por el punto exterior $P(1,-3)$.\recorda{$P$ no es de la curva: escribe la tangente en un punto genérico $(a,a^2)$ e impón que pase por $P$.}

\ex{2} Halla los parámetros en cada caso:\begin{tasks}(1)
  \task \pau{Madrid, extraordinaria 2022} $a$ y $b$ para que $f(x)=ax+\dfrac bx$ verifique $f(2)=4$ y $f'(2)=0$.
  \task \pau{Murcia, septiembre 2017} $a$ y $b$ para que $f(x)=\dfrac{ax+b}{x^2+1}$ cumpla $f(0)=1$ y $f'(0)=1$.
  \task \pau{Andalucía, ordinaria 2024} $a$ y $b$ para que sea horizontal la recta tangente a la gráfica de $h(x)=x^3+ax^2+3x+b$ en el punto $P(1,2)$.
\end{tasks}

\ex{2} \pau{La Rioja, ordinaria 2024} Escribe la recta tangente a $y=\dfrac1x$ en el punto $\left(3,\frac13\right)$ y comprueba que el segmento de esta recta comprendido entre los ejes queda dividido en dos partes iguales por el punto de tangencia.

\ex{2} \pau{Castilla-La Mancha, extraordinaria 2021} Calcula las rectas tangente y normal a la gráfica de $f(x)=\dfrac{2x^2+2x-2}{3x^2+3}$ en el punto de abscisa $x=1$.

\ex{2} \pau{Andalucía, extraordinaria 2024} Halla $a$, $b$ y $c$ sabiendo que la gráfica de $f(x)=a+b\cos x+c\sin x$ tiene por recta tangente $y=1$ en el punto de abscisa $x=\frac\pi2$, y que la recta $y=x-1$ corta la gráfica en el punto de abscisa $x=0$.

\ex{3} \pau{La Rioja, extraordinaria 2024} Escribe, si existen, las rectas tangentes a la curva $f(x)=|x|\,e^{-x}$ en los puntos de abscisa $x=0$ y $x=-1$.\recorda{Escribe primero la función a trozos: $|x|=-x$ si $x<0$.}

\ex{2} \pau{Aragón, ordinaria 2022} Para la función $f(x)=\begin{cases}\sqrt{-x}+e^3&\text{si }x\leq0\\[3pt](1-x)^{a/x}&\text{si }x>0\end{cases}$ (que es continua para $a=-3$), calcula la recta tangente en $x=-4$.


% =====================================================
\seccio{3}{Reglas de derivación y derivadas sucesivas}

\ex{1} Deriva y simplifica:\begin{tasks}(2)
  \task $f(x)=x^3\sqrt{x}-\dfrac{2}{\sqrt[3]{x}}$
  \task $f(x)=x\,e^x\sin x$
  \task $f(x)=\dfrac{x^2-1}{x^2+1}$
  \task $f(x)=\dfrac{x^3}{x-1}$
  \task $f(x)=\dfrac{1-\cos x}{\sin x}$
  \task $f(x)=e^x\left(x^2-2x+2\right)$
\end{tasks}\recorda{Escribe raíces e inversas como potencias. Producto de tres: $(fgh)'=f'gh+fg'h+fgh'$. Cociente: $\left(\frac fg\right)'=\frac{f'g-fg'}{g^2}$.}

\ex{1} \pau{Cantabria, ordinaria 2024} Considera la función $f(x)=x\ln x$ con $x>0$. Calcula la derivada de $f(x)$.

\ex{2} Comprueba que $F(x)=x\arctan x-\frac12\ln(x^2+1)$ tiene por derivada $\arctan x$.

\ex{2} Derivadas sucesivas:\begin{tasks}(1)
  \task Calcula $f'''(x)$ si $f(x)=x^4-2x^3+x$.
  \task Halla la derivada $n$-ésima de $e^{-2x}$, de $x\,e^x$, de $\dfrac{1}{x+1}$ y de $\ln x$.
  \task ¿Cuál es la derivada 2027.ª de $\cos x$?
\end{tasks}\recorda{Calcula tres o cuatro, busca el patrón y escríbelo en función de $n$. Las derivadas de $\sin x$ y $\cos x$ se repiten cada cuatro.}

\ex{2} Halla $a$, $b$ y $c$ para que $f(x)=ax^2+bx+c$ cumpla $f(1)=2$, $f'(1)=3$ y $f''(1)=4$.

\ex{3} Comprueba que la función $y=e^{-x}\sin x$ verifica la ecuación $y''+2y'+2y=0$.


% =====================================================
\seccio{4}{Regla de la cadena}

\ex{1} Deriva estas potencias, raíces y exponenciales:\begin{tasks}(3)
  \task $f(x)=(3x^2-1)^4$
  \task $f(x)=\sqrt{x^2-2x}$
  \task $f(x)=\dfrac{1}{(x^2+1)^3}$
  \task $f(x)=\sqrt[3]{2x+5}$
  \task $f(x)=e^{3x-1}$
  \task $f(x)=2^{x^2}$
\end{tasks}\recorda{$\left[f(g(x))\right]'=f'(g(x))\cdot g'(x)$: derivamos la de fuera, dejando dentro $g(x)$, y multiplicamos por la derivada de la de dentro.}

\ex{2} Deriva estas funciones trigonométricas y trigonométricas inversas:\begin{tasks}(3)
  \task $f(x)=\sin^2(3x)$
  \task $f(x)=\cos(x^2)$
  \task $f(x)=\tan^3x$
  \task $f(x)=\arctan(2x)$
  \task $f(x)=\arcsin\sqrt x$
  \task $f(x)=\arccos e^{-x}$
\end{tasks}

\ex{2} Deriva y simplifica (aprovecha las propiedades de los logaritmos antes de derivar):\begin{tasks}(3)
  \task $f(x)=\ln(x^2-3x)$
  \task $f(x)=\ln\sqrt{\dfrac{1+x}{1-x}}$
  \task $f(x)=\log_2(x^3+1)$
  \task $f(x)=\ln(\cos x)$
  \task $f(x)=\ln\left(x+\sqrt{x^2+1}\right)$
  \task $f(x)=\ln\left(\tan\frac x2\right)$
\end{tasks}

\ex{2} Deriva y simplifica (hay más de una regla en cada función):\begin{tasks}(2)
  \task $f(x)=x\,e^{-x^2}$
  \task $f(x)=(x-1)^2(x+2)^3$
  \task $f(x)=\sqrt{x+\sqrt x}$
  \task $f(x)=e^{\sin x}\cos x$
  \task $f(x)=\sin(\ln x)+\cos(\ln x)$
\end{tasks}

\ex{2} \pau{Murcia, septiembre 2020} Calcula las derivadas de las funciones siguientes:\begin{tasks}(2)
  \task $f(x)=(x^2-2)\ln x$
  \task $g(x)=e^{4x^3+2}$
\end{tasks}

\ex{2} \pau{Murcia, junio 2018} Halla las derivadas de las funciones siguientes:\begin{tasks}(2)
  \task $g(x)=x^2\ln(x^2)$
  \task $h(x)=e^{-3x+x^2}$
\end{tasks}

\ex{2} \pau{Andalucía, ordinaria 2021} Calcula la derivada de las funciones siguientes:\begin{tasks}(2)
  \task $f(x)=\ln\left(\dfrac{x-1}{x+1}\right)$
  \task $g(x)=x^3\cdot e^{2x^2}$
\end{tasks}

\ex{2} \pau{Andalucía, ordinaria (suplente) 2025} Calcula la derivada de las funciones siguientes:\begin{tasks}(2)
  \task $g(x)=\ln\left(\dfrac{x^2-1}{x^2+1}\right)$
  \task $h(x)=(2x-1)\,e^{x^2-x}$
\end{tasks}

\ex{3} \pau{Andalucía, ordinaria 2024} Calcula la derivada de las funciones siguientes:\begin{tasks}(2)
  \task $f(x)=(x^2+2)^3\cdot e^{-2x}$
  \task $g(x)=\dfrac{\ln(1-x^3)}{(1-2x^2)^2}$
\end{tasks}

\ex{2} \pau{Murcia, septiembre 2018} Halla $a$ y $b$ para que $f(x)=x^3\ln(2x+5)+ax+b$ cumpla $f(0)=2$ y $f'(0)=1$.


% =====================================================
\seccio{5}{Derivación logarítmica}

\ex{2} Deriva aplicando la derivación logarítmica:\begin{tasks}(2)
  \task $y=x^x$
  \task $y=x^{\sin x}$
  \task $y=x^{\ln x}$
  \task $y=x^{1/x}$
  \task $y=(x^2+1)^x$
  \task $y=(\ln x)^x\ \ (x>1)$
  \task $y=(\cos x)^{\sin x}\ \ \left(0<x<\frac\pi2\right)$
  \task $y=x^{x^2}$
\end{tasks}\recorda{Tomamos logaritmos: $\ln y=v\ln u$. Derivamos los dos miembros: $\dfrac{y'}{y}=v'\ln u+v\dfrac{u'}{u}$. Despejamos $y'$ y sustituimos $y$.}

\ex{2} Deriva con la derivación logarítmica (es más rápido que con las reglas del producto y del cociente):\begin{tasks}(2)
  \task $y=\dfrac{x^3(x+1)^5}{\sqrt{x^2+4}}\ \ (x>0)$
  \task $y=\dfrac{(x-1)^2\sqrt[3]{x+2}}{(x+3)^4}\ \ (x>1)$
\end{tasks}

\ex{2} Considera la función $f(x)=x^x$, con $x>0$.\begin{tasks}(1)
  \task Halla las rectas tangentes a la gráfica en $x=1$ y en $x=2$.
  \task ¿En qué punto la tangente es horizontal?
  \task Explica por qué es incorrecto derivarla como si fuera una potencia, $(x^x)'=x\cdot x^{x-1}$, o como si fuera una exponencial, $(x^x)'=x^x\ln x$.
\end{tasks}


% =====================================================
\seccio{6}{Derivación implícita}

\ex{2} Calcula $y'$ en cada caso:\begin{tasks}(2)
  \task $x^2+y^2=4x$
  \task $x^3+y^3=3xy$
  \task $x^2y+xy^2=6$
  \task $e^{xy}=x+y$
  \task $\sqrt x+\sqrt y=3$
  \task $x=y+\ln y$
\end{tasks}\recorda{Derivamos los dos miembros respecto a $x$ pensando que $y=y(x)$: $(y^2)'=2y\,y'$, \ $(xy)'=y+x\,y'$, \ $(e^{y})'=e^{y}\,y'$. Después despejamos $y'$.}

\ex{2} Halla la recta tangente (y, en el primer caso, también la normal) a cada curva en el punto indicado. Comprueba primero que el punto es de la curva.\begin{tasks}(1)
  \task $x^2+xy+y^2=7$ en $(2,1)$
  \task $x^3+y^3=9xy$ en $(2,4)$
  \task $\dfrac{x^2}{16}+\dfrac{y^2}{9}=1$ en $\left(\frac{16}{5},\frac95\right)$
  \task $y^2=x^3$ en $(1,1)$
\end{tasks}

\ex{3} ¿En qué puntos de la curva $x^2+xy+y^2=3$ la recta tangente es horizontal? ¿Y vertical?

\ex{3} Calcula la derivada segunda $y''$ y exprésala en función de $x$ y $y$ (o solo de $y$):\begin{tasks}(3)
  \task $x^2+y^2=r^2$
  \task $x^2-y^2=1$
  \task $y^3+y=x$
\end{tasks}

\ex{2} Deduce, derivando implícitamente, las derivadas de $y=\ln x$ (a partir de $e^y=x$), de $y=\arcsin x$ (a partir de $\sin y=x$) y de $y=\arctan x$ (a partir de $\tan y=x$).

\ex{3} Las dos técnicas a la vez. El punto $(2,4)$ es de la curva $x^y=y^x$, ya que $2^4=4^2$. Halla la pendiente de la tangente en este punto.\recorda{Toma primero logaritmos ($y\ln x=x\ln y$) y después deriva implícitamente.}


% =====================================================
\seccio{7}{Continuidad y derivabilidad}

\ex{2} Calcula las derivadas laterales en el punto indicado y di si la función es derivable en él:\begin{tasks}(2)
  \task $f(x)=|x+1|$ en $x=-1$
  \task $f(x)=|x^2-9|$ en $x=3$
  \task $f(x)=x\,|x-1|$ en $x=1$
  \task $f(x)=\sqrt[3]{x^2}$ en $x=0$
  \task $f(x)=|\sin x|$ en $x=\pi$
  \task $f(x)=x^2|x|$ en $x=0$
\end{tasks}\recorda{$f'(a^\pm)=\displaystyle\lim_{h\to0^\pm}\frac{f(a+h)-f(a)}{h}$. Es derivable si ambas existen, son finitas y son iguales.}

\ex{1} Razona si son verdaderas o falsas. Si son falsas, pon un contraejemplo:\begin{tasks}(1)
  \task Si $f$ es continua en $x=a$, entonces es derivable en $x=a$.
  \task Si $f$ no es continua en $x=a$, entonces no es derivable en $x=a$.
  \task Si $f$ y $g$ no son derivables en $x=a$, $f+g$ tampoco lo es.
  \task Si $f$ es derivable en $x=a$, $|f|$ también lo es.
\end{tasks}

\ex{2} \pau{Galicia, ordinaria 2026} Dibuja, cuando sea posible, la gráfica de una función continua pero no derivable en $x=0$.

\ex{3} Sea $f(x)=x^2\sin\dfrac1x$ si $x\neq0$ y $f(0)=0$. Demuestra, aplicando la definición, que $f$ es derivable en $x=0$. ¿Pasa lo mismo con $g(x)=x\sin\dfrac1x$, $g(0)=0$?\recorda{Si una función tiende a 0 y la otra está acotada, el producto tiende a 0.}


% =====================================================
\seccio{8}{Derivabilidad de funciones a trozos}

\ex{2} Estudia la continuidad y la derivabilidad de las funciones siguientes:\begin{tasks}(2)
  \task $f(x)=\begin{cases}e^x&\text{si }x<0\\[3pt]x+1&\text{si }x\geq0\end{cases}$
  \task $g(x)=\begin{cases}3-x^2&\text{si }x<1\\[3pt]\dfrac2x&\text{si }x\geq1\end{cases}$
\end{tasks}\recorda{Primero la continuidad. Después derivamos trozo a trozo (sin el igual en los puntos de cambio) y, en cada punto de cambio donde es continua, comparamos $f'(a^-)$ y $f'(a^+)$.}

\ex{2} \pau{Andalucía, extraordinaria 2024} Estudia la continuidad y la derivabilidad de\[f(x)=\begin{cases}-x^2+4x+3&\text{si }x<4\\[3pt]2x-5&\text{si }x\geq4\end{cases}\]

\ex{2} \pau{Andalucía, ordinaria (suplente) 2025} Estudia la continuidad y la derivabilidad de $f$ en el punto de abscisa $x=-2$:\[f(x)=\begin{cases}10+\dfrac{5x}{2}&\text{si }x\leq-2\\[3pt]x^2+1&\text{si }-2<x<2\\[3pt]10-\dfrac{5x}{2}&\text{si }x\geq2\end{cases}\]

\ex{2} \pau{Andalucía, ordinaria 2024} La velocidad del viento, en km/h, prevista para cierto día en una zona de Sierra Nevada viene dada por\[v(t)=\begin{cases}t^2-8t+60&\text{si }0\leq t\leq10\\[3pt]-t^2+32t-140&\text{si }10<t\leq24\end{cases}\]donde $t$ es el tiempo en horas. Comprueba que la función $v$ es continua y derivable.

\ex{2} \pau{Navarra, ordinaria 2024} Estudia la continuidad (clasificando las discontinuidades) y la derivabilidad de\[f(x)=\begin{cases}-x^2-4x&\text{si }x<-1\\[3pt]4-x&\text{si }-1\leq x<1\\[3pt]x^2-6x+8&\text{si }x\geq1\end{cases}\]

\ex{2} \pau{Andalucía, extraordinaria 2021} Estudia la continuidad y la derivabilidad de\[f(x)=\begin{cases}2^{x+1}&\text{si }x<0\\[3pt]x^2-2x&\text{si }x\geq0\end{cases}\]

\ex{3} \pau{Madrid, extraordinaria 2022 (ampliado)} Estudia la continuidad y la derivabilidad en $x=0$ de\[f(x)=\begin{cases}x&\text{si }x\leq0\\[3pt]x\ln x&\text{si }x>0\end{cases}\]

\ex{2} \pau{Asturias, ordinaria 2025 (ampliado)} Estudia la continuidad y la derivabilidad de\[f(x)=\begin{cases}\dfrac x2&\text{si }0\leq x<4\\[3pt]3-(x-5)^2&\text{si }4\leq x\end{cases}\]

\ex{3} \pau{Andalucía, ordinaria 2022 (ampliado)} Considera la función continua\[f(x)=\begin{cases}\dfrac1x&\text{si }x<-1\\[3pt]ax+b&\text{si }-1\leq x<1\\[3pt]\dfrac{x^2}{x+1}&\text{si }x\geq1\end{cases}\]\begin{tasks}(1)
  \task Calcula $a$ y $b$.
  \task Estudia la derivabilidad de $f$.
\end{tasks}

\ex{2} Estudia la derivabilidad de\begin{tasks}(2)
  \task $f(x)=\begin{cases}\sin x&\text{si }x\leq0\\[3pt]\ln(1+x)&\text{si }x>0\end{cases}$
  \task $g(x)=\begin{cases}\cos x&\text{si }x\leq0\\[3pt]1+x^2&\text{si }x>0\end{cases}$
\end{tasks}

\ex{3} \pau{Andalucía, extraordinaria 2021} El número de diagnosticados por PCR en Andalucía, en miles de personas, se aproxima por\[f(t)=\begin{cases}-t^2+2t-0{,}3&\text{si }0{,}2\leq t\leq1{,}8\\[3pt]0{,}1t-0{,}12&\text{si }1{,}8<t\leq5\\[3pt]-0{,}5t^2+8{,}3t-28{,}62&\text{si }5<t\leq10\end{cases}\]donde $t$ es el tiempo en meses. Estudia la continuidad y la derivabilidad de $f$.

\ex{3} Escribe a trozos y estudia la derivabilidad de:\begin{tasks}(2)
  \task $f(x)=|x-1|+|x+1|$
  \task $g(x)=x\,|x-2|$
  \task $h(x)=e^{-|x|}$
  \task $k(x)=|x^2-5x+4|$
\end{tasks}


% =====================================================
\seccio{9}{Funciones a trozos con parámetros}

\ex{2} \pau{Andalucía, ordinaria 2022} Determina $a$ y $b$ para que $f$ sea continua y derivable en todo su dominio:\[f(x)=\begin{cases}6x-3&\text{si }x\leq1\\[3pt]ax^2+bx+2&\text{si }x>1\end{cases}\]\recorda{En cada punto de cambio: continuidad ($\lim^-=\lim^+=f(a)$) y derivabilidad ($f'(a^-)=f'(a^+)$). Con dos parámetros, un sistema de dos ecuaciones.}

\ex{2} \pau{Andalucía, extraordinaria 2020} Calcula $a$ y $b$ para que la función sea continua y derivable en su dominio:\[f(x)=\begin{cases}2+\dfrac{a}{x-1}&\text{si }x<0\\[3pt]a+be^x&\text{si }x\geq0\end{cases}\]

\ex{2} Calcula los valores de $a$ para que $f$ sea continua y derivable:\begin{tasks}(1)
  \task \pau{Islas Baleares, ordinaria 2022} \ $f(x)=\begin{cases}x^2-3x+2&\text{si }x<0\\[3pt]x^2+ax+2&\text{si }x\geq0\end{cases}$
  \task \pau{Islas Baleares, ordinaria 2021} \ $f(x)=\begin{cases}x^2-3x+2&\text{si }x<0\\[3pt]e^{ax}+1&\text{si }x\geq0\end{cases}$
\end{tasks}

\ex{2} \pau{Islas Baleares, extraordinaria 2021} Considera\[f(x)=\begin{cases}-4x+a&\text{si }x\leq-2\\[3pt]x^2-5&\text{si }-2<x<1\\[3pt]-7x+3&\text{si }x\geq1\end{cases}\]\begin{tasks}(1)
  \task Calcula los valores de $a$ para que $f$ sea continua.
  \task ¿Es $f$ derivable para $a=1$?
\end{tasks}

\ex{3} \pau{Aragón, ordinaria 2022} Siendo $a$, $b$ parámetros reales, se considera\[f(x)=\begin{cases}x^2+2\sqrt3&\text{si }x\leq0\\[3pt]\sqrt{ax+b}&\text{si }0<x\leq3\\[3pt]\dfrac72-\dfrac x6&\text{si }x>3\end{cases}\]\begin{tasks}(1)
  \task Determina el valor de los parámetros para que $f$ sea continua.
  \task Para estos valores, analiza si $f$ es derivable en $x=0$ y en $x=3$.
\end{tasks}

\ex{2} \pau{Andalucía, ordinaria (suplente) 2025} El nivel de concentración de un alumno durante un examen de 5 horas viene dado por\[f(t)=\begin{cases}-t^2+2t+10&\text{si }0\leq t\leq2{,}5\\[3pt]t^2+at+b&\text{si }2{,}5<t\leq5\end{cases}\]donde $t$ es el tiempo en horas. ¿Con qué nivel empieza el examen? Determina $a$ y $b$ para que $f$ sea continua y derivable en $t=2{,}5$.

\ex{3} \pau{Cataluña, ordinaria 2026 (ampliado)} Determina los valores de $m$ que hacen que $f$ sea continua en todo su dominio y, para estos valores, estudia su derivabilidad:\[f(x)=\begin{cases}5e^{2x}&\text{si }x\leq0\\[3pt](x+m)^2+1&\text{si }0<x<2\\[3pt]1&\text{si }2\leq x\end{cases}\]

\ex{2} \pau{Extremadura, extraordinaria 2023 (ampliado)} Halla $a$ y $b$ para que $f$ sea continua en $x=1$ y su gráfica pase por $(-1,5)$. ¿Es derivable en $x=1$?\[f(x)=\begin{cases}2x^2+ax+b&\text{si }x\leq1\\[3pt]\ln x&\text{si }x>1\end{cases}\]

\ex{2} \pau{Aragón, extraordinaria 2021 (ampliado)} Calcula los valores de $a\neq0$ para que $f$ sea continua y estudia si, para alguno de estos valores, es derivable en $x=1$:\[f(x)=\begin{cases}5-ax^2&\text{si }x\leq1\\[3pt]\dfrac{6}{ax}&\text{si }x>1\end{cases}\]

\ex{2} Halla los parámetros para que cada función sea derivable en su dominio:\begin{tasks}(2)
  \task $f(x)=\begin{cases}ax^2+bx&\text{si }x\leq1\\[3pt]\ln x&\text{si }x>1\end{cases}$
  \task $g(x)=\begin{cases}ax^2+b&\text{si }x<1\\[3pt]\dfrac1x&\text{si }x\geq1\end{cases}$
  \task $h(x)=\begin{cases}e^{ax}+b&\text{si }x\leq0\\[3pt]\sin2x+3&\text{si }x>0\end{cases}$
  \task $k(x)=\begin{cases}a+bx&\text{si }x\leq0\\[3pt]\arctan x&\text{si }x>0\end{cases}$
\end{tasks}

\ex{2} Halla los parámetros en cada caso:\begin{tasks}(1)
  \task \pau{Galicia, ordinaria 2025} ¿Qué condición tienen que cumplir $k$ y $m$ para que $f(x)=\begin{cases}kx^2+2x&\text{si }x\leq1\\[3pt]x^2-m&\text{si }x>1\end{cases}$ sea continua en $x=1$? ¿Y para qué valores es derivable en $x=1$?
  \task \pau{Galicia, extraordinaria 2024} Si $f(x)=\begin{cases}x^2+bx-1&\text{si }x\leq0\\[3pt]\dfrac{k-xe^x}{x}&\text{si }x>0\end{cases}$, ¿qué valor de $k$ la hace continua en $x=0$ para cualquier $b$? ¿Para qué $b$ y $k$ es derivable en $x=0$?
  \task \pau{Canarias, junio 2019} Halla $a$ y $b$ para que $f(x)=\begin{cases}a-x&\text{si }x\leq1\\[3pt]\dfrac bx-\ln x&\text{si }x>1\end{cases}$ sea derivable en todo su dominio y escribe la función resultante.
\end{tasks}\recorda{Si un trozo tiene la forma $\frac{\dots}{x}$ y tiene que ser continuo en 0, el numerador se tiene que anular en 0 (si no, el límite es infinito).}

\ex{2} Halla los parámetros para que cada función sea continua y derivable en el punto de cambio:\begin{tasks}(1)
  \task \pau{Castilla-La Mancha, ordinaria 2021} $f(x)=\begin{cases}x^2-ax+1&\text{si }x<0\\[3pt]be^x&\text{si }x\geq0\end{cases}$
  \task \pau{Castilla-La Mancha, ordinaria 2026} $f(x)=\begin{cases}\dfrac{x^2-1}{x-1}&\text{si }x<1\\[3pt]ax+b&\text{si }x\geq1\end{cases}$
  \task \pau{Andalucía, junio 2018} $k\neq0$ para que sea derivable $f(x)=\begin{cases}3-kx^2&\text{si }x\leq1\\[3pt]\dfrac{2}{kx}&\text{si }x>1\end{cases}$
\end{tasks}

\ex{3} \pau{Madrid, ordinaria 2026} Sea $f(x)=\begin{cases}3x^2-\sin(ax)+b&\text{si }x<0\\[3pt]\dfrac{2e^x}{1+x^2}&\text{si }x\geq0\end{cases}$.\begin{tasks}(1)
  \task Calcula $a$ y $b$ para que sea continua y derivable en $x=0$.
  \task Halla la recta tangente a la gráfica en $x=2$.
\end{tasks}

\vfill
\begin{center}\color{muted}\small\faCheckCircle[regular]\ \ ¡Ánimo! Comprueba tus derivadas con la calculadora de los apuntes interactivos.\end{center}

\end{document}
